% Source printed pp. 457-466 / PDF pp. 465-474.
\section{半稳定向量丛的模叠}\label{src:9.2}
\begin{epigraph}
于是可以看出, ``扭曲层''理论所带来的成功主要在心理层面: 将模叠视为自身独立的对象 (而不是要压缩成概形的东西), 更容易保留那些可以重新流回簇的世界的信息.
\epigraphauthor{Max Lieblich \cite{Lie18}}
\end{epigraph}
我们讨论半稳定向量丛模叠 $\cBun_{r,d}^{\mathrm{ss}}$ 的性质. 重点包括半稳定性的开放性 (\bookref{9.2.6}), 半稳定性的有界性 (\bookref{9.2.8}), 以及稳定轨迹 $\cBun_{r,d}^{\mathrm s}$ 非空 (定理~\bookref{9.2.16}). 最后一项性质利用扩张叠 (\S\bookref{9.2.4-subsection}) 归纳证明. 我们还说明 $\cBun_{r,d}^{\mathrm s}$ 是代数空间 $M_{r,d}^{\mathrm s}$ 上带 $\mathbb G_m$-带的 gerbe (推论~\bookref{9.2.26}). 最后引入固定行列式的向量丛叠 (\S\bookref{9.2.7-subsection}).
\subsection{曲线上向量丛叠的回顾}\label{src:9.2.1-subsection}
设 $C$ 是代数闭域 $\Bbbk$ 上亏格 $g\geq2$ 的光滑连通射影曲线. 我们已经建立了参数化 $C$ 上凝聚层的叠 $\Coh:=\Coh(C)$ 的若干性质, 现在予以回顾. 首先, $\Coh$ 是代数叠 (定理~\bookref{4.1.23}), 局部有限型 (例~\bookref{4.3.3}), 且对角仿射 (例~\bookref{4.3.15}). 定理~\bookref{4.1.23} 还蕴含参数化秩为 $r$, 次数为 $d$ 的凝聚层的子叠 $\Coh_{r,d}\subseteq\Coh$ 既开又闭, 而向量丛的子叠 $\cBun\subseteq\Coh$ 和 $\cBun_{r,d}\subseteq\Coh_{r,d}$ 是开子叠. $\Coh_{r,d}$ 的每个拟紧开子叠都同构于某个 $N$ 对应的商叠 $[Q_N'/\mathrm{GL}_{P(N)}]$, 其中 $Q_N'\subset\Quot^P(\mathcal O_C(-N)^{\oplus P(N)}/C/\Bbbk)$ 是射影 Quot 概形的开子概形, $P$ 是任意秩为 $r$, 次数为 $d$ 的凝聚层的 Hilbert 多项式.

利用形变理论 (\bookref{C.2.11}), 我们在命题~\bookref{4.7.9} 中证明 $\Coh_{r,d}$ 在 $\Spec\Bbbk$ 上光滑, 维数为 $r^2(g-1)$. 事实上, 若 $[F]\in\Coh_{r,d}(\Bbbk)$, 则 $T_{\Coh_{r,d},[F]}\cong\operatorname{Ext}^1(F,F)$, Riemann--Roch 定理计算出其维数为 $r^2(g-1)$ (见习题~\bookref{9.1.4}). 叠 $\Coh_{r,d}$ 和 $\cBun_{r,d}$ 满足普遍闭性的赋值判据 (习题~\bookref{4.8.14}), 但不拟紧 (习题~\bookref{9.2.3}).

叠 $\Coh_{r,d}$ 和 $\cBun_{r,d}$ 不是 Deligne--Mumford 叠, 因为 $\underline{\operatorname{Aut}}(F)$ 总包含一个 $\mathbb G_m$ (只要 $(r,d)\ne(0,0)$). 所以 $\Coh_{r,d}$ 和 $\cBun_{r,d}$ 不分离: 若它们分离, 则每个自同构群都固有且仿射, 因而有限. 另一方面, 已知 $\Coh_{r,d}$ 是 $\Theta$-完备且 S-完备的 (命题~\bookref{7.9.30}), 这对在 \S\bookref{9.3.2-subsection} 建立良模空间的存在性很有用.
\begin{exercise}\label{src:9.2.1}
设 $F$ 是 $C$ 上秩为 $r$, 次数为 $d$ 的凝聚层, $N$ 是使得 $\mathrm H^1(C,F(N))=0$ 且 $F(N)$ 由整体截面生成的整数, 即 $F$ 可表为商 $[\mathcal O_C(-N)\otimes V\twoheadrightarrow F]\in Q_N'$, 其中 $V\cong\mathrm H^0(C,F(N))$ 是维数 $P(N)$ 的向量空间. 证明存在正合列, 其中各项有如下所示的等同:
\[
\begin{aligned}
0\to\underbrace{\Hom_{\mathcal O_C}(F,F)}_{\operatorname{End}(F)}
\to\underbrace{\Hom_{\mathcal O_C}(\mathcal O_C(-N)\otimes V,F)}_{\operatorname{End}(V)}
\to\underbrace{\Hom_{\mathcal O_C}(K,F)}_{T_{Q_N',[\mathcal O_C(-N)\otimes V\to F]}}\\
\to\underbrace{\operatorname{Ext}_{\mathcal O_C}^1(F,F)}_{T_{\Coh_{r,d},[F]}}
\to\underbrace{\operatorname{Ext}_{\mathcal O_C}^1(\mathcal O_C(-N)\otimes V,F)}_0.
\end{aligned}
\]
\end{exercise}
\begin{remark}[叠 \texorpdfstring{$\cBun_{r,d}(\mathbb P^1)$}{Bun(P1)}]\label{src:9.2.2}
虽然 Birkhoff--Grothendieck 定理 (\bookref{9.1.5}) 简单地分类了 $\mathbb P^1$ 上向量丛的同构类, 向量丛叠 $\cBun_{r,d}(\mathbb P^1)$ 却更复杂. $\cBun_{r,d}(\mathbb P^1)$ 并不是按向量丛同构类分解的不交并. 例如, 考虑由 $\operatorname{Ext}_{\mathcal O_{\mathbb P^1}}^1(\mathcal O(1),\mathcal O(-1))\cong\Bbbk$ 的非零元素定义的短正合列 $0\to\mathcal O(-1)\to\mathcal O^{\oplus2}\to\mathcal O(1)\to0$. 万有扩张定义 $\mathbb G_m$-等变映射 $\mathbb A^1\to\cBun_{2,0}(\mathbb P^1)$, 将 $1$ 映到 $[\mathcal O^{\oplus2}]$, 将 $0$ 映到 $[\mathcal O(-1)\oplus\mathcal O(1)]$. 等价地, 可将它视为映射 $[\mathbb A^1/\mathbb G_m]\to\cBun_{2,0}(\mathbb P^1)$; 其存在性也来自命题~\bookref{7.9.6} 对来自 $[\mathbb A^1/\mathbb G_m]$ 的映射的刻画. 这蕴含 $[\mathcal O(-1)\oplus\mathcal O(1)]\in|\cBun_{2,0}(\mathbb P^1)|$ 位于 $[\mathcal O^{\oplus2}]$ 的闭包中. 此外, 因为 $\cBun_{2,0}(\mathbb P^1)$ 包含每个 $d>0$ 对应的 $\mathcal O(d)\oplus\mathcal O(-d)$, $\cBun_{2,0}(\mathbb P^1)$ 不有界, 即不是有限型的.
\end{remark}
\begin{exercise}\label{src:9.2.3}
证明代数叠 $\cBun_{r,d}(C)$ 不有界.
\end{exercise}
\subsection{半稳定性的开放性}\label{src:9.2.2-subsection}
我们证明稳定性与半稳定性是开放性质, 并推出叠 $\cBun_{r,d}^{\mathrm s}$ 和 $\cBun_{r,d}^{\mathrm{ss}}$ 是 $\cBun_{r,d}$ 的开子叠.
\begin{theorem}[稳定性/半稳定性的开放性]\label{src:9.2.4}
设 $C$ 是代数闭域 $\Bbbk$ 上亏格 $g\geq2$ 的光滑连通射影曲线. 设 $S$ 是 $\Bbbk$-概形, $F$ 是 $C_S$ 上秩为 $r$, 次数为 $d$ 的向量丛族. 使 $F_s:=F|_{C_s}$ 半稳定的点 $s\in S$ 组成开轨迹. 类似地, 使 $F_{\overline s}:=F|_{C\times\Spec\kappa(\overline s)}$ 稳定的点 $s\in S$ 组成开轨迹.
\end{theorem}
\begin{remark}\label{src:9.2.5}
以上两种表述的差别来自这样一个事实: 与半稳定性不同, 稳定性在域扩张下不总保持 (习题~\bookref{9.1.30}).
\end{remark}
\begin{proof}
由 \S\bookref{B.3} 的极限方法, 可假设 $S$ 在 $\Bbbk$ 上有限型. 还可假设 $S$ 连通. 我们将使 $F_s$ 不半稳定的点 $s\in S$ 的轨迹表示为有限个固有态射的像之并. 设 $P(t):=rt+(d+r(1-g))$ 为秩 $r$, 次数 $d$ 的向量丛的 Hilbert 多项式. 由习题~\bookref{9.1.28}, 存在整数 $D$, 使每个纤维的每个子层 $E\subseteq F_s$ 满足 $\deg E\leq D$. 由定义, $F_s$ 不半稳定, 当且仅当存在子层 $E\subseteq F_s$, 满足 $\mu(F)>d/r$. 因为 $\operatorname{rk}E\leq r$ 且 $\deg E\leq D$, $F_s$ 的不稳定化子层 $E$ 的 $(\operatorname{rk}E,\deg E)$ 只有有限种可能. 对每个可能的秩次数 $(r',d')$, 令 $P'(t)=r't+(d'+r'(1-g))$ 为秩 $r'$, 次数 $d'$ 的向量丛的 Hilbert 多项式. 令 $Q(t)=P(t)-P'(t)$ 为商 $F_s/E$ 的 Hilbert 多项式. 由 Quot 概形的射影性 (\bookref{2.1.3}), 存在参数化 Hilbert 多项式为 $Q$ 的 $F$ 的商的概形 $\Quot^Q(F/C_S/S)$, 它在 $S$ 上射影. $\Quot^Q(F/C_S/S)\to S$ 的像闭, 并与使 $F_s$ 有秩 $r'$, 次数 $d'$ 的不稳定化子层的点 $s\in S$ 的轨迹等同. 所以 $S$ 的非半稳定轨迹闭.

类似地, 若 $F_{\overline s}$ 半稳定但不稳定, 令 $h=\gcd(r,d)$, 则对某个整数 $0<k<\gcd(r,d)$, 存在秩 $kr/h$, 次数 $kd/h$ 的子层 $0\ne E\subsetneq F_{\overline s}$. 上述论证推广后说明使 $F_{\overline s}$ 不稳定的点 $s\in S$ 的轨迹闭.
\end{proof}
\begin{corollary}\label{src:9.2.6}
叠 $\cBun_{r,d}^{\mathrm s}$ 和 $\cBun_{r,d}^{\mathrm{ss}}$ 是 $\cBun_{r,d}$ 的开子叠; 特别地, 是 $\Bbbk$ 上局部有限型代数叠.
\end{corollary}
\subsection{半稳定性的有界性}\label{src:9.2.3-subsection}
曲线上半稳定向量丛的有界性的关键是下述引理.
\begin{lemma}\label{src:9.2.7}
若 $E$ 是 $C$ 上半稳定向量丛且 $\mu(E)>2g-1$, 则 $E$ 由整体截面生成, 且 $\mathrm H^1(C,E)=0$.
\end{lemma}
\begin{proof}
由 Serre 对偶 (\bookref{6.1.3}), $\mathrm H^1(C,E)\cong\Hom_C(E,\Omega_C)^\vee$. 线丛 $\Omega_C$ 半稳定, 斜率为 $2g-2$. 因为 $E$ 也半稳定且 $\mu(E)>2g-2$, 由引理~\bookref{9.1.15}(1), $\Hom_C(E,\Omega_C)=0$. 所以 $\mathrm H^1(C,E)=0$. 相同的论证说明对任意点 $p\in C(\Bbbk)$, $\mathrm H^1(C,E\otimes\mathcal O_C(-p))=0$, 因为 $E\otimes\mathcal O_C(-p)$ 是半稳定向量丛 (引理~\bookref{9.1.13}(4)), 斜率为 $\mu(E)-1>2g-2$. 将此消失结论应用于短正合列 $0\to E\otimes\mathcal O_C(-p)\to E\to E\otimes\kappa(p)\to0$, 得到对每个 $p\in C$, $\mathrm H^0(C,E)\to E\otimes\kappa(p)$ 满射.
\end{proof}
\begin{theorem}[半稳定性的有界性]\label{src:9.2.8}
设 $C$ 是代数闭域 $\Bbbk$ 上亏格 $g\geq2$ 的光滑连通射影曲线. 对每对整数 $r>0$ 和 $d$, 代数叠 $\cBun_{r,d}^{\mathrm{ss}}$ 在 $\Bbbk$ 上有限型.
\end{theorem}
\begin{proof}
选取点 $p\in C(\Bbbk)$ 及整数 $N$, 使 $d/r+N>2g-1$. 对每个秩为 $r$, 次数为 $d$ 的半稳定向量丛 $E$, 向量丛 $E\otimes\mathcal O_C(Np)$ 由引理~\bookref{9.1.13}(4) 半稳定, 且斜率为 $\mu(E\otimes\mathcal O_C(Np))=\mu(E)+N>2g-1$. 由引理~\bookref{9.2.7}, $E\otimes\mathcal O_C(Np)$ 由整体截面生成, 且 $\mathrm H^1(C,E\otimes\mathcal O_C(Np))=0$. 所以 $E\otimes\mathcal O_C(Np)$ 可表为 $V\otimes\mathcal O_C$ 的商, 其中 $V$ 是固定维数 $\mathrm h^0(C,E\otimes\mathcal O_C(Np))=\chi(C,E\otimes\mathcal O_C(Np))$ (与 $E$ 无关) 的向量空间, 因而 $E$ 可表为商
\[
V\otimes\mathcal O_C(-Np)\to E.
\]
令 $P(t)=rt+(d+r(1-g))$, 由定理~\bookref{2.1.3}, Quot 概形 $\Quot^P(V\otimes\mathcal O_C(-Np)/C)$ 在 $\Bbbk$ 上有限型. 由半稳定性的开放性 (\bookref{9.2.4}), 存在参数化半稳定商的开子概形 $Q'\subseteq\Quot^P(V\otimes\mathcal O_C(-Np)/C)$. $Q'$ 上的万有商定义满射态射 $Q'\to\cBun_{r,d}^{\mathrm{ss}}$ (事实上, 将 $\cBun_{r,d}^{\mathrm{ss}}$ 表为商叠 $[Q'/\mathrm{GL}_{P(N)}]$). 因为 $Q'$ 拟紧, $\cBun_{r,d}^{\mathrm{ss}}$ 也拟紧.
\end{proof}
\begin{exercise}[后文需要]\label{src:9.2.9}
固定整数 $r,d$ 及实数 $\gamma$.
\begin{enumerate}[(a)]
\item 证明由秩为 $r$, 次数为 $d$ 的半稳定向量丛 $F$ 及满足 $\mu_{\min}(E)\geq\gamma$ 的凝聚子层 $E\subseteq F$ 组成的对的轨迹有界. (回顾 $\mu_{\min}(E)$ 表示 $E$ 的 Harder--Narasimhan 滤过 (\bookref{9.1.21}) 因子的最小斜率.)
\item 对每个 $\gamma$, 证明使所有凝聚子层 $E\subseteq F$ 满足 $\mu(E)\leq\mu(F)+\gamma$ 的秩 $r$, 次数 $d$ 向量丛 $F$ 的轨迹有界, 即该轨迹包含于 $\cBun_{r,d}$ 的某个有限型子叠中.
\end{enumerate}
\end{exercise}
